\chapter{Conclusiones}

El índice horario es una herramienta indispensable para la correcta configuración 
de sistemas eléctricos de potencia. A través de este parámetro, es posible 
determinar no solo la relación de fase entre los distintos niveles de tensión, 
sino también asegurar la compatibilidad entre equipos para su operación conjunta.

Se concluye que:
\begin{itemize}
    \item El grupo de conexión más versátil para la distribución en baja tensión 
    es el \textbf{Dy11}, debido a su capacidad para manejar cargas desequilibradas 
    y filtrar armónicos.
    \item La verificación del índice horario es un paso crítico previo a la puesta 
    en servicio de cualquier transformador que deba operar en paralelo con otros, 
    siendo la causa más común de fallas graves en subestaciones si se omite.
    \item La representación del reloj simplifica enormemente la interpretación de 
    un fenómeno complejo como es el desfase electromagnético en máquinas 
    rotativas y estáticas.
\end{itemize}

El cumplimiento de los estándares de conexión garantiza la estabilidad de la red 
y la protección de los activos de la infraestructura eléctrica.
